\subsection{Heron}
\paragraph{TODO 1} verify topics via experiment

\subsubsection{Hardware}
There's MiR mobile platform and a UR5e arm. The MiR platform is differential drive.
It has some Lidar and ultrasonic distance sensors. By default, the arm and the base
can not move together, but there is a specific secret hack to circumvent this.

\subsubsection{Software}
There is a whole ROS2 stack, but some parts require additional work. 
By default you control the robot over WiFi, which comes from an onboard router.
The onboard computer supports putting a whole Docker image on it,
meaning one can easily get the best possible control performance.

Despite the ROS stack, SMC arm control relies on ur\_rtde, just as when there's just the manipulator, 
but calls to ur\_rtde are now in callbacks instead.
This choice has been made as ur\_rtde works perfectly fine, and there is
absolutely no benefit in ROS beyond localization (and the utility of this
for whole-body control testing is already very debatable).

The base is controllable only via ROS though.
Base is controlled with \fbox{\texttt{/cmd\_vel}} which runs at 30Hz.

\subsubsection{Localization}
Look at localization subsection in the navigation section.
Odometry is at \fbox{\texttt{/odometry}} and runs at 60Hz.
There's AMCL for localization, but \textbf{carefully consider} whether you really need it.
If you are simply testing a whole-body controller to open a door for example,
there is absolutely no need for localization whatsoever --- just use odometry.
